\documentclass[11pt,a4paper]{article}
\usepackage[utf8]{inputenc}
\usepackage[T1]{fontenc}
\usepackage{lmodern}
\usepackage{amsmath,amssymb,amsthm}
\usepackage{mathtools}
\usepackage[ngerman]{babel}
\usepackage{booktabs}
\usepackage{microtype}
\usepackage{geometry}
\usepackage{hyperref}

\geometry{a4paper, top=2.2cm, bottom=2.2cm, left=2.5cm, right=2.5cm}
\hypersetup{colorlinks=true, linkcolor=black, citecolor=black, urlcolor=black}

\newtheorem{theorem}{Satz}[section]
\newtheorem{lemma}[theorem]{Lemma}
\newtheorem{proposition}[theorem]{Proposition}
\theoremstyle{definition}
\newtheorem{definition}[theorem]{Definition}
\theoremstyle{remark}
\newtheorem{remark}[theorem]{Bemerkung}

\numberwithin{equation}{section}

\title{\textbf{\texorpdfstring{Collatz im $E_8$-Rahmen:\\
Einbettung, Drift-Heuristik und die offene L\"ucke}{Collatz im E8-Rahmen: Einbettung, Drift-Heuristik und die offene Luecke}}}
\author{\textbf{Thomas Hoffbauer}\thanks{Bamberger Modell f\"ur Mathematische Physik, Zahlentheorie und Spektralgeometrie.
Korrespondenz: \texttt{t.hoffbauer@posteo.de}}}
\date{14.~September 2026}

\begin{document}

\maketitle

\begin{abstract}
\noindent
Die Collatz-Vermutung besagt, dass jede positive ganze Zahl unter der Syracuse-Abbildung
$T$ den trivialen Zyklus \(\{1,2\}\) erreicht.
Zwei Hindernisse stehen einem Beweis entgegen: unbeschr\"ankte Trajektorien und fremde Zyklen.
Diese Note bettet \(T\) als \emph{Analogie} in die \(2\)-adische Vervollst\"andigung und in die
Spektralgeometrie von \(E_8\) ein.
Die logarithmische Drift \(\mathbb{E}[\Delta\log n]=\tfrac12\log(3/4)<0\) und die
Hoffbauer-Klammer \(\sqrt{D_r^2-D_t^2}\) sind Heuristiken bzw.\ Sonden aus dem EABC-Modell;
sie liefern keine Uniformschranke an \(T^k(n)\).
Die klassische Zyklusgleichung bleibt diophantisch.
Die Kongruenz \(2^K\equiv 3^M\pmod{24}\) ist daf\"ur \emph{nicht} notwendig;
Dedekind-Stufe \(24\) und \(\Delta_{41}=861\) ersetzen sie nicht.
Die Collatz-Vermutung bleibt unbewiesen.
\end{abstract}

\noindent\textbf{MSC (2020):} 11B83, 37P35, 11S82, 37A25 \\
\textbf{Schl\"usselw\"orter:} Collatz-Vermutung (offen), Syracuse-Abbildung, \(2\)-adische Dynamik,
\(E_8\), \(\mathrm{Spin}(8)\)-Trialit\"at, Hoffbauer-Klammer, Drift-Heuristik.

\section{Problemstellung}

Auf \(\mathbb{N}=\{1,2,3,\dots\}\) sei die komprimierte Syracuse-Abbildung
\begin{equation}
T(n)=\begin{cases}
n/2, & n\equiv 0\pmod{2},\\[0.4em]
(3n+1)/2, & n\equiv 1\pmod{2}.
\end{cases}
\end{equation}
Die Collatz-Vermutung fordert: zu jedem \(n\in\mathbb{N}\) existiert ein \(k\in\mathbb{N}_0\) mit \(T^k(n)=1\).
Danach gilt der triviale \(2\)-Zyklus \(1\mapsto 2\mapsto 1\).
In der unkomprimierten Form \(n\mapsto 3n+1\) (ungerade) bzw.\ \(n\mapsto n/2\) (gerade)
ist derselbe Orbit der \(3\)-Zyklus \(\{1,4,2\}\).

Trotz Verifikation weit \"uber \(n>2^{68}\) und trotz Taos Satz, dass \emph{fast alle} Orbits
(logarithmische Dichte) unter jedes langsam wachsende \(f(n)\to\infty\) fallen~\cite{Tao2019},
fehlt ein Beweis f\"ur \emph{jedes} \(n\).
Die Schwierigkeit hat zwei klassische Formen~\cite{Lagarias1985}:
\begin{enumerate}
\item \textbf{Divergenz.} Es k\"onnte ein \(n_0\) geben mit \(\limsup_{k\to\infty}T^k(n_0)=\infty\).
\item \textbf{Fremde Zyklen.} Es k\"onnte eine Periode \(L\ge 3\) in der Syracuse-Z\"ahlung geben,
die nicht \(\{1,2\}\) ist.
Steiner schloss \(1\)-Zyklen au\ss{}er dem trivialen aus~\cite{Steiner1977};
Simons und de Weger erweiterten das auf kleine \(m\)-Zyklen~\cite{SimonsDeWeger2005}.
Ein genereller Ausschluss fehlt.
\end{enumerate}

Eine beschr\"ankte Trajektorie in \(\mathbb{N}\) trifft nach dem Schubfachprinzip einen Zyklus.
W\"are zus\"atzlich gezeigt, dass (i)~jede Trajektorie beschr\"ankt ist und (ii)~nur \(\{1,2\}\)
als Zyklus vorkommt, so folgte Collatz.
Genau (i) und (ii) sind offen.
Diese Note trennt Analogie, Heuristik und L\"ucke; sie behauptet keinen Hauptsatz.

\section{\texorpdfstring{Analogie: \(\mathbb{Z}_2\) und \(E_8\)}{Analogie: Z2 und E8}}

Die Fallunterscheidung \(n\bmod 2\) legt die \(2\)-adische Vervollst\"andigung nahe.
Im Ring \(\mathbb{Z}_2=\varprojlim\mathbb{Z}/2^k\mathbb{Z}\) besitzt jedes Element eine
eindeutige Entwicklung \(n=\sum_{j\ge 0}a_j 2^j\) mit \(a_j\in\{0,1\}\).
Die Abbildung \(T\) setzt sich dort zu einer stetigen Transformation fort
(Bernstein--Lagarias, \(2\)-adische Collatz-Dynamik).
Maßtreue bez\"uglich des Haar-Maßes auf \(\mathbb{Z}_2\) ist ein Satz \"uber den
\emph{kompakten} Ring, kein Transfer auf den positiven Tr\"ager \(\mathbb{N}\subset\mathbb{Z}\).
Ergodizit\"at auf \(\mathbb{Z}_2\) schließt divergierende positive Orbits nicht aus.

Im EABC-Rahmen lesen wir zwei weitere Analogien, nicht \"Aquivalenzen von dynamischen Systemen:
\begin{itemize}
\item Division durch \(2\) als Verschiebung entlang des geraden Untergitters \(D_8\subset E_8\).
\item Der Faktor \(3\) als Ordnung von \(\mathrm{Out}(\mathrm{Spin}(8))\cong\mathfrak{S}_3\).
\end{itemize}
Die Fano-Ebene der Oktonionen und eine Drehung um \(2\pi/3\) auf \(S^7\) sind geometrische
Bilder derselben Zahl \(3\).
Sie ersetzen weder die Fallunterscheidung in \(T\) noch eine H\"ohe auf \(\mathbb{N}\).

\section{Drift-Heuristik}

F\"ur unabh\"angige Parit\"aten mit Wahrscheinlichkeit \(1/2\) gilt die klassische Erwartung
\begin{equation}
\mathbb{E}[\Delta\log n]
=\tfrac12\log\bigl(\tfrac12\bigr)+\tfrac12\log\bigl(\tfrac32\bigr)
=\tfrac12\log\bigl(\tfrac34\bigr)<0.
\end{equation}
Numerisch ist \(\log(3/4)\approx -0{,}28768\).
Das ist die Standard-Heuristik f\"ur \emph{typische} Orbits, kein Ausschluss pathologischer
Parit\"atsfolgen, in denen ungerade Schritte dauerhaft dominieren.
Die kritische Dichte ungerader Schritte f\"ur Wachstum ist
\(\alpha_{\mathrm{crit}}=\log 2/\log 3\approx 0{,}6309\)
in der unkomprimierten Z\"ahlung; unabh\"angige M\"unzw\"urfe liegen darunter, abh\"angige
Folgen m\"ussen das nicht.

Tao~\cite{Tao2019} macht daraus einen Satz \"uber logarithmische Dichte.
Einzeltrajektorien bleiben außerhalb der Reichweite maßtheoretischer Mittelwerte.

\section{Hoffbauer-Klammer als Sonde, nicht als Schranke}

Im EABC-Modell ist die hyperbolische Defektklammer
\begin{equation}
\Delta_{\mathrm{EABC}}
=[D_r,D_t]_{\mathcal{H}}
:=\sqrt{D_r^2-D_t^2},
\end{equation}
mit den festgehaltenen rationalen Daten
\begin{equation}
D_r=\frac{8}{137\,\varphi},\qquad
D_t=\frac{126}{49457},\qquad
\varphi=\frac{1+\sqrt{5}}{2},\qquad
126=\frac{\tau(3)}{2}.
\end{equation}
Hier ist \(\tau(3)=252\) der Ramanujan-Wert der Spitzenform \(\Delta\in S_{12}(\mathrm{SL}_2(\mathbb{Z}))\),
und \(49457=361\cdot 137\) mit der Hecke-Stufe \(c=361/126\).
Numerisch gilt \(\Delta_{\mathrm{EABC}}\approx 0{,}03599954\); das ist der Nachkommaanteil
der konstruierten Größe \(\alpha^{-1}_{\mathrm{EABC}}=137+\Delta_{\mathrm{EABC}}\)~\cite{EABC2026}.

\begin{remark}[Keine Uniformschranke]
Eine unbeschr\"ankte Collatz-Trajektorie w\"urde weder \(\tau(3)\) noch \(D_t>0\) aufheben.
Die Planck-Sonde \(\ell_P=2\sqrt{\alpha^{-1}_{\mathrm{EABC}}}\approx 23{,}412\) des geometrischen
Einheitensystems \(e=c=1\), \(G=4\) ist dimensionslos im EABC-Rahmen; sie ist keine Schranke
an \(T^k(n)\in\mathbb{N}\).
Dissipation in \(\mathbb{H}^\perp\subset\mathbb{O}\) ist ein Bild, kein Lyapunov-Funktional
auf \(\mathbb{N}\).
Insbesondere folgt \emph{nicht}
\[
\sup_{k\ge 0}T^k(n)<\infty\qquad\text{f\"ur alle }n\in\mathbb{N}.
\]
\end{remark}

Die Klammer bleibt eine sinnvolle \emph{Probe}: dieselbe algebraische \(3\), die in
\(\mathrm{Out}(\mathrm{Spin}(8))\) und in \(T_3\) auftritt, steuert den ungeraden Collatz-Schritt.
Das ist Motiv, kein Abstiegssatz.

\section{Zyklen: diophantische Gleichung}

Ein Zyklus mit \(M\) ungeraden Schritten und insgesamt \(K\) Divisionen durch \(2\)
(unkomprimierte Z\"ahlung) erf\"ullt eine Gleichung der Form
\begin{equation}
\label{eq:cycle}
(2^K-3^M)\,n_0
=\sum_{j=0}^{M-1} 3^{M-1-j}\,2^{A_j},
\end{equation}
mit strikt wachsenden Exponenten \(0=A_0<A_1<\dots<A_{M-1}<K\).
F\"ur positives \(n_0\) braucht man \(2^K>3^M\), also \(K>M\log 3/\log 2\).
Das ist klassisch~\cite{Lagarias1985}.
Notwendig ist, dass \(2^K-3^M\) die rechte Seite \emph{teilt}.
Nicht notwendig ist
\begin{equation}
\label{eq:false-cong}
2^K\equiv 3^M\pmod{24}.
\end{equation}

Die Kongruenz~\eqref{eq:false-cong} ist sogar f\"ur \(K\ge 3\) und \(M\ge 1\)
stets falsch: \(3^M\bmod 24\in\{3,9\}\), w\"ahrend \(2^K\bmod 24\in\{8,16\}\)
(nie \(0\), weil \(2^K\) nie durch \(3\) teilbar ist).
Die leere Schnittmenge \(\{3,9\}\cap\{8,16\}=\emptyset\) ist daher eine Tatsache
\"uber Potenzreste, keine Bedingung an~\eqref{eq:cycle}.
Der triviale Zyklus hat in der unkomprimierten Z\"ahlung \(K=2\), \(M=1\)
und umgeht~\eqref{eq:false-cong} nur deshalb, weil \(K<3\);
das rettet~\eqref{eq:false-cong} nicht als Kriterium f\"ur \(K\ge 3\).

\begin{remark}[Modularit\"at als Mnemotechnik]
Die Zahl \(24\) tritt als \(8\cdot 3=\mathrm{rk}(E_8)\cdot 3\),
als Wurzelzahl von \(D_4\) und als Gewicht von \(\Delta(\tau)\) auf.
Das sind parallele Vorkommen derselben ganzen Zahl, keine Homologiebedingung
an Collatz-Zyklen in \(H_1(T^8/\mathcal{H}^\times,\mathbb{Z})\).
Ebenso ist \(\lfloor h\rfloor=861=\Delta_{41}\) im EABC-Planck-Bild eine
Zirkulationssonde, keine teilerfremde Obstruktion gegen Linearkombinationen
von Potenzen von \(2\) und \(3\).
\end{remark}

Was bleibt, ist die echte Literatur: Ausschluss kleiner \(m\)-Zyklen,
explizite untere Schranken an die Zyklusl\"ange, und die offene diophantische
Frage, ob~\eqref{eq:cycle} außer dem trivialen Orbit positive L\"osungen hat.

\section{Was eine Schließung br\"auchte}

Drei Schichten sind zu trennen.
\begin{itemize}
\item \textbf{Bestand.}
Syracuse-Abbildung \(T\), Zyklusgleichung~\eqref{eq:cycle},
Steiner, Simons--de Weger, Tao (fast alle Orbits).
\item \textbf{Analogie / Sonde.}
Fortsetzung von \(T\) auf \(\mathbb{Z}_2\);
\(\times 3\leftrightarrow\mathfrak{S}_3\); \(/2\leftrightarrow D_8\);
Hoffbauer-Klammer; \(\tau(3)=252\); \(\Delta_{41}\).
\item \textbf{Offen.}
Uniformschranke an \(T^k(n)\);
Ausschluss aller fremden Zyklen;
Transfer vom kompakten Ring \(\mathbb{Z}_2\) auf \(\mathbb{N}_{>0}\).
\end{itemize}

\vspace{0.6em}
Ein Beweis m\"usste mindestens eines der folgenden liefern:
\begin{enumerate}
\item ein Lyapunov-Funktional \(V:\mathbb{N}\to\mathbb{R}\) mit \(V(T(n))\le V(n)-\delta\)
außerhalb einer expliziten endlichen Menge, oder eine vergleichbare H\"ohenfunktion
mit kontrolliertem Anstieg;
\item eine vollst\"andige Restklassen-Partition (Zylinder-Cover), in der jeder Zweig
nach endlich vielen Schritten unter eine feste Schranke f\"allt, mit Lean-Claim
\(\mathtt{True}\) statt Sample-Coverage;
\item eine diophantische Absch\"atzung, die~\eqref{eq:cycle} f\"ur alle \(K\ge 3\)
außer dem trivialen Orbit ausschließt, ohne eine nicht notwendige Kongruenz.
\end{enumerate}
Finite Sample-Cover, \(2\)-adische Ergodizit\"at und EABC-Sonden erf\"ullen keines
dieser drei Kriterien~\cite{CollatzFast2026}.
Das Residuen-Pruning \(\rho D>B\) (Gatter~D) ist ein deterministisches Branch-and-Bound auf \(\mathbb{N}\),
t\"otet aber nur kontrahierende Zylinder \(D>0\); expandierende W\"orter wie \((2,1)^{\infty}\) bleiben.

\section{Schluss}

Die Collatz-Dynamik l\"asst sich im EABC-Rahmen als dissipative Bewegung auf einem
diskreten Spinor \emph{lesen}.
Die ungerade \(3\) und die Hecke-Stufe \(T_3\) sind dieselbe algebraische Zahl in
zwei Sprachen.
Das rechtfertigt die Einbettung als Forschungsprogramm.

Es rechtfertigt nicht den Satz, jede Trajektorie erreiche \(1\).
Divergenz und fremde Zyklen sind weiterhin die zwei Bastionen.
Die Hoffbauer-Klammer und die Dedekind-Stufe \(24\) sind Sonden an \(\alpha^{-1}\),
keine Schranken an \(T\).
Die Collatz-Vermutung bleibt offen.

\begin{thebibliography}{99}
\bibitem{Tao2019}
T.~Tao,
\emph{Almost all orbits of the Collatz map attain almost bounded values},
arXiv:1909.03562 [math.PR] (2019).

\bibitem{Steiner1977}
R.~P.~Steiner,
\emph{A theorem on the Syracuse problem},
Proc.\ 7th Manitoba Conf.\ Numerical Math.\ (1977), 553--559.

\bibitem{SimonsDeWeger2005}
J.~Simons, B.~de Weger,
\emph{Theoretical and computational bounds for \(m\)-cycles of the \(3x+1\)-problem},
Acta Arith.\ \textbf{117} (2005), 51--70.

\bibitem{Lagarias1985}
J.~C.~Lagarias,
\emph{The \(3x+1\) problem and its generalizations},
Amer.\ Math.\ Monthly \textbf{92} (1985), 3--23.

\bibitem{EABC2026}
T.~Hoffbauer,
\emph{Hyperbolische Abschirmung und Dedekind-Ordnung:
ein \(E_8\)-Zugang zum CODATA-Fenster von \(\alpha^{-1}\)},
Preprint, Bamberg (2026).

\bibitem{CollatzFast2026}
Bamberg / Kepler--Hurwitz,
\emph{Adaptive Cover und Zylinder-Partitionen -- Abschlussbericht},
eabc-renorm, 25.~Juli 2026.
\end{thebibliography}

\end{document}
